\documentclass[10pt,a4paper]{amsart}
\usepackage[margin=19mm]{geometry}
\usepackage{amsmath,amssymb,mathtools}
\usepackage[scr=rsfs,cal=euler]{mathalfa}
\usepackage{xfrac}
\usepackage[unicode,hidelinks,bookmarks=false]{hyperref}
\newcommand{\ie}{i.e.\ }
\newcommand{\cf}{cf.\ }
\newcommand{\resp}{respectively\ }
\newcommand{\Subsection}{\subsection}
\providecommand{\nobreakdash}{\nobreak\hbox{-}}
\setlength{\emergencystretch}{2em}
\multlinegap0pt
\newcommand{\C}{\ensuremath{\mathbb{C}}}
\usepackage{amssymb}
\usepackage{bm}
\usepackage{enumerate}
\usepackage[normalem]{ulem}
\newtheorem{theorem}{Theorem}
\title{Energy asymptotics of holomorphic functions with application to Calder\'{o}n-Zygmund theory in $\C$}
\author{Hongrong Chen, Guokuan Shao, Jujie Wu, Wei Xia}
\date{Source: arXiv:2609.02601v1 (CC BY 4.0)}
\begin{document}
\maketitle

\noindent\emph{Source and licence.} Energy asymptotics of holomorphic functions with application to Calder\'{o}n-Zygmund theory in $\C$. Version arXiv:2609.02601v1 (CC BY 4.0), distributed by arXiv under the Creative Commons Attribution 4.0 International licence, http://creativecommons.org/licenses/by/4.0/, as recorded in the arXiv licence metadata for this exact version and shown on the abstract page https://arxiv.org/abs/2609.02601v1. Full licence notice: \texttt{licenses/arxiv-2609.02601-CC-BY-4.0.txt}.\par\smallskip

\noindent\textit{Editorial scope.} Counted unit: the article's theorem Linfty (source0.tex lines 376--403). The article's complete original proof of it is delivered with it (source0.tex lines 2323--2485, one \texttt{proof} environment of 163 lines, which the article places away from the statement). No mathematics is written, repaired or re-derived: the statement and the proof are the article's own lines, in the article's own notation, and no retained line is substituted or re-flowed.  No further source-local context is needed: every cross-reference of every retained line resolves inside the two delivered spans.  Nothing was omitted from any retained span, so the package carries no omission record.  No retained line contains a citation, so the wrapper carries no bibliography and no bibliography entry.  No retained line contains a figure environment, an image-inclusion command or drawing code, so no image is delivered and the package declares no asset dependency.  Editorial formatting only, in addition: an AMS \texttt{amsart} preamble that restates the article's own declarations and macros used by the retained lines, namely the environments \texttt{theorem} on separate counters, together with the standard packages those lines need.  The retained environments print in this document as Theorem 1 in order of appearance, whereas the article prints the same items with the numbering of its own full document.  The complete native source of the article as distributed in the arXiv e-print for this exact version is bundled unchanged as \texttt{sources/209216/source0.tex}.
\begin{theorem}\label{Linfty}
Let $U \subset \mathbb{C}$ be a bounded domain, and let $f \in \mathcal{O}(\overline{U})$ be a nonconstant holomorphic function such that $Z(f) \cap U = \{a_1, \dots, a_N\}$ is non-empty ($N \ge 1$). Define the polynomial
\[
P(z) := \prod_{j=1}^N (z-a_j),
\]
and set
\[
u(z) :=
\begin{cases}
P(z)^2 \log |f(z)|^2, & z \in U \setminus Z(f),\\[4pt]
0, & z \in Z(f) \cap U.
\end{cases}
\]
Then the following statements hold:
\begin{enumerate}
    \item $u \in C^1(U)$;
    \item $g \in L^\infty_{\mathrm{loc}}(U)$ and $\Delta u = g$ in the sense of distributions on $U$, where
    \begin{equation}
    g(z) :=
    \begin{cases}
    \Delta u(z), & z \in U \setminus Z(f),\\[4pt]
    0, & z \in Z(f) \cap U;
    \end{cases}
    \end{equation}
    \item $u \notin W^{2,\infty}_{\mathrm{loc}}(U)$.
\end{enumerate}
Consequently, $u$ serves as a one-dimensional counterexample to Calder\'on--Zygmund estimates at $p = \infty$.
\end{theorem}

\begin{proof}[Proof of Theorem \ref{Linfty}]

%\medskip
%\noindent\textbf{Proof of (1).}
We prove the three assertions in order.

$(1)$ Fix $a_j\in Z(f)\cap U$, and let $m_j\ge1$ be the multiplicity of $a_j$ as a zero of $f$.
In a small neighborhood of $a_j$, we may write
$$
f(z)=(z-a_j)^{m_j}h_j(z),
\qquad h_j(a_j)\neq0.
$$
Since $P$ has a simple zero at $a_j$, we may also write
$$
P(z)=(z-a_j)Q_j(z),
\qquad Q_j(a_j)\neq0,
$$
where $Q_j$ is holomorphic near $a_j$. Hence $P(z)^2=(z-a_j)^2Q_j(z)^2$.

Near $a_j$, we have
$$
P(z)^2\log|f(z)|^2
=
(z-a_j)^2Q_j(z)^2
\left(
m_j\log|z-a_j|^2+\log|h_j(z)|^2
\right).
$$
The function $(z-a_j)^2\log|z-a_j|^2$ is $C^1$ near $a_j$. Therefore
$P(z)^2\log|f(z)|^2$ extends as a $C^1$ function across $a_j$.
Since this argument applies to every zero $a_j\in Z(f)\cap U$, and $u$ is smooth on
$U\setminus Z(f)$, we obtain $u\in C^1(U)$.

%\medskip
%\noindent\textbf{Proof of (2).}
$(2)$ On $U\setminus Z(f)$, since $P^2$ is holomorphic, we have
$$
\frac{\partial u}{\partial\bar z}
=
P(z)^2\frac{\overline{f'(z)}}{\overline{f(z)}}.
$$
Differentiating with respect to $z$, and using that
$\overline{f'(z)}/\overline{f(z)}$ is anti-holomorphic on $U\setminus Z(f)$, we get
$$
\frac{\partial^2 u}{\partial z\partial\bar z}
=
2P(z)P'(z)\frac{\overline{f'(z)}}{\overline{f(z)}}.
$$
Therefore, on $U\setminus Z(f)$,
$$
\Delta u
=
8P(z)P'(z)\frac{\overline{f'(z)}}{\overline{f(z)}}.
$$
Thus, $g$ is exactly the classical Laplacian of $u$ on $U\setminus Z(f)$.

We claim that $g\in L^\infty_{\mathrm{loc}}(U)$. It suffices to check boundedness near each zero
$a_j$. Near $a_j$, write $P(z)=(z-a_j)Q_j(z)$ with $Q_j(a_j)\neq0$. Then
$P(z)P'(z)=O(|z-a_j|)$. Moreover,
$$
\frac{\overline{f'(z)}}{\overline{f(z)}}
= \frac{m_j}{\overline{z - a_j}} + \frac{\overline{h'_j(z)}}{\overline{h_j(z)}}.
$$
Thus,
$$
P(z)P'(z)\frac{\overline{f'(z)}}{\overline{f(z)}}=O(1)
$$
near $a_j$. Away from $Z(f)\cap U$, the function $g$ is smooth. Hence,
$g\in L^\infty_{\mathrm{loc}}(U)$.

It remains to prove that $\Delta u=g$ holds on $U$ in the sense of distributions. Since
the statement is local, it is enough to check that no extra Dirac mass appears at each zero.
Fix $a_j$, and choose $r>0$ such that $D(a_j,r)\Subset U$ and $D(a_j,r)$ contains no zero of
$f$ other than $a_j$. Let $\eta\in C_c^\infty(D(a_j,r))$. For $0<\varepsilon<r$, set
$\Omega_\varepsilon:=D(a_j,r)\setminus\overline{D(a_j,\varepsilon)}$. Since $u$ is smooth on
$\Omega_\varepsilon$ and satisfies $\Delta u=g$ there, Green's identity gives
$$
\int_{\Omega_\varepsilon}u\Delta\eta\,dA
-
\int_{\Omega_\varepsilon}g\eta\,dA
=
\int_{|z-a_j|=\varepsilon}
\left(
u\frac{\partial\eta}{\partial\nu}
-
\eta\frac{\partial u}{\partial\nu}
\right)ds.
$$
From the expression in the proof of (1), as $z\to a_j$ we have
$u(z)=O(|z-a_j|^2|\log|z-a_j||)$ and
$|\nabla u(z)|=O(|z-a_j||\log|z-a_j||)$. Hence, on $|z-a_j|=\varepsilon$,
$$
|u(z)|\le C\varepsilon^2|\log\varepsilon|,
\qquad
|\nabla u(z)|\le C\varepsilon|\log\varepsilon|.
$$
Since the length of the circle $|z-a_j|=\varepsilon$ is $2\pi\varepsilon$, the boundary integral
is bounded by
$$
C\varepsilon^3|\log\varepsilon|+C\varepsilon^2|\log\varepsilon|,
$$
which tends to $0$ as $\varepsilon\to0^+$. Letting $\varepsilon\to0^+$ gives
$$
\int_{D(a_j,r)}u\Delta\eta\,dA
=
\int_{D(a_j,r)}g\eta\,dA.
$$
Thus $\Delta u=g$ distributionally across $a_j$. Since $a_j$ was arbitrary, $\Delta u=g$ holds
on $U$ in the sense of distributions.

%\medskip
%\noindent\textbf{Proof of (3).}
$(3)$
Fix $a_j\in Z(f)\cap U$, and write $\xi=z-a_j$. As above, near $a_j$ we have
$f(z)=\xi^{m_j}h_j(z)$ with $h_j(a_j)\neq0$, and $P(z)=\xi Q_j(z)$ with
$Q_j(a_j)\neq0$. Set $R_j(z):=Q_j(z)^2$. Then $P(z)^2=\xi^2R_j(z)$ and
$R_j(a_j)\neq0$.

On a punctured neighborhood of $a_j$, a direct computation gives
\begin{align}
\frac{\partial^2 u}{\partial z^2}
&=
(P^2)''(z)\log|f(z)|^2
+
2(P^2)'(z)\frac{f'(z)}{f(z)}
+
P(z)^2
\left(
\frac{f''(z)}{f(z)}
-
\left(\frac{f'(z)}{f(z)}\right)^2
\right).
\label{eq:second-derivative-u-theorem13}
\end{align}
The last two terms in \eqref{eq:second-derivative-u-theorem13} are bounded near $a_j$.
Indeed, $(P^2)'(z)=O(|\xi|)$ and $f'(z)/f(z)=m_j/\xi+O(1)$, while $P(z)^2=O(|\xi|^2)$ and
$$
\frac{f''(z)}{f(z)}
-
\left(\frac{f'(z)}{f(z)}\right)^2
=
-\frac{m_j}{\xi^2}+O(1).
$$
On the other hand, $(P^2)''(a_j)=2R_j(a_j)\neq0$, and
$$
\log|f(z)|^2
=
m_j\log|\xi|^2+O(1).
$$
Therefore,
\begin{align}
\frac{\partial^2 u}{\partial z^2}
=
2m_jR_j(a_j)\log|z-a_j|^2+O(1)
\qquad \text{as } z\to a_j.
\label{eq:singular-second-derivative-u-theorem13}
\end{align}
Since $R_j(a_j)\neq0$, the asymptotic formula
\eqref{eq:singular-second-derivative-u-theorem13} shows that $\partial_z^2u$ is essentially unbounded in every neighborhood of $a_j$. Hence, $\partial_z^2u\notin L^\infty_{\mathrm{loc}}(V)$ for any open neighborhood
$V\subset U$ of $a_j$.
Therefore,
$u\notin W^{2,\infty}_{\mathrm{loc}}(U)$. This proves (3).
\end{proof}

\end{document}
